\section{Initial Benchmark Experiments on \sc{ConflictQA}}
\label{sec:experiment}

We study pluralistic groundedness as an evidence aggregation problem: given a set
of retrieved documents that support competing positions, does a model produce
an answer that faithfully reflects the available evidence? We evaluate this
question using ConflictQA \citep{wan-etal-2024-evidence}, a benchmark designed
to measure how language models reason over conflicting evidence.

\subsection{ConflictQA}

ConflictQA consists of contested yes-or-no questions paired with evidence
passages that support opposing answers. The questions span nineteen domains across the social 
and natural sciences, humanities, and applied fields. Each example contains a query, a set of
passages supporting a \textit{yes} answer, and a set of passages supporting a
\textit{no} answer. For example, a question such as \textit{Is aspartame linked
to cancer?} is paired with evidence representing both the claim that such a link
exists and the claim that current evidence does not support such a link.

\subsection{Evidence Aggregation Probes}

We construct controlled probes that vary properties of the retrieved evidence
while keeping the underlying question fixed. These probes test whether models
respond to evidence itself or to incidental features of its presentation.

\paragraph{Expertise and Rhetorical Authority.}
We next examine how models weight signals of authority when aggregating
conflicting evidence. Authority can be a legitimate component of evidence
aggregation: for some questions, domain expertise should influence how
evidence is weighted. For example, when evaluating whether aspartame is linked
to cancer, the judgment of a qualified scientist may deserve more weight than
the opinion of a non-expert. However, models may also mistake rhetorical
confidence or formal style for expertise.

To separate these effects, we create evidence variants that modify the
presentation style of passages while holding their underlying claims constant.
The authoritative condition uses formal, confident language, while the casual
condition uses conversational language. This probe tests whether models respond
to authority cues appropriately.


\paragraph{Evidence Balance.}
We next vary the number of passages supporting each answer. We denote an
evidence configuration as $y$ vs. $n$, where $y$ is the number of passages
supporting a yes answer and $n$ is the number supporting a no answer. We
evaluate six configurations:

\begin{itemize}
\item $3v3$: balanced evidence with three passages supporting each answer;
\item $2v2$: balanced evidence with two passages supporting each answer;
\item $2v0$: evidence supporting yes only;
\item $0v2$: evidence supporting no only;
\item $4v2$: asymmetric evidence favoring yes;
\item $2v4$: asymmetric evidence favoring no.
\end{itemize}

The balanced conditions test whether models produce consistent judgments when
the amount of evidence changes but the balance of perspectives remains fixed.
The asymmetric conditions test whether models track differences in evidence
coverage, while the one-sided conditions test whether models introduce
unsupported disagreement when only one position is represented.

Together, these probes test whether models aggregate evidence according to
appropriate signals rather than relying on simple heuristics. A faithful
aggregator should respond to meaningful differences between evidence sources,
such as expertise when domain knowledge is relevant, while avoiding
unwarranted sensitivity to superficial features such as rhetorical confidence.
It should also adjust its conclusions when the distribution of evidence changes
without introducing unsupported perspectives.
